\section{Audio Transformers：Raw Waveform 与多尺度表示}

第三篇工作不再预测下一个 sample，而是从一秒 raw waveform 识别 sound events。它借鉴 ViT 的核心思路：把输入切成 patches，经 learned front end 映射为 embeddings，再用 Transformer 聚合；同时用 pooling 或 wavelet 让后层看到更宽时间范围。这里的“raw”仍包含前端切块与线性投影，并非把 16,000 个 samples 原封不动全部做 global attention。

\lecturefigure{slide-38-audio-transformers-title.jpg}{第三篇工作：无卷积的 large-scale raw-audio understanding。}{Stanford Online 官方视频 00:37:50--00:38:21；Audio Transformers v1, arXiv:2105.00335v1。}

\subsection{FSD50K 任务与公平比较边界}

课堂使用 FSD50K：约五万条 human-labeled sound events、约 200 类。Slide 中的课堂设置约为 40,000 training clips 与 10,000 validation/test clips，总量约百小时；输入下采样为 16 kHz，并以一秒 waveform 为基本样本。目标是 multi-label audio content classification，而不是 waveform reconstruction。

\lecturefigure{slide-39-fsd50k-dataset-setup.jpg}{FSD50K 与实验设置：raw waveform、固定时长、multi-label content prediction。}{Stanford Online 官方视频 00:37:54--00:38:42。}

对样本 $x$ 与 $C$ 个标签，模型输出 logits $s_c$，常用 binary cross-entropy：
\[
\mathcal{L}_{\mathrm{BCE}}
=-\sum_{c=1}^{C}
\left[y_c\log\sigma(s_c)+(1-y_c)\log(1-\sigma(s_c))\right].
\]
其中 $y_c\in\{0,1\}$ 表示第 $c$ 类是否出现。评价使用 mean average precision（mAP），它综合每类 precision--recall ranking，再对类别平均；因此不能和前文的 top-5 next-sample accuracy 直接比较。

\begin{warningbox}{同叫 Accuracy，任务完全不同}
前文 raw synthesis 在 256 个 sample states 中做 single-label next-step prediction；这里 FSD50K 是 multi-label event detection，报告 mAP。二者数值大小没有可比性，不能说 0.537 “低于” 85\%。
\end{warningbox}

\subsection{为什么在 Intermediate Embedding 上做 Multi-Scale}

Raw waveform 包含从瞬态到事件的多尺度结构。单纯堆叠等长 Transformer layers，虽然每层理论上可全局寻址，但表示维度和 token rate 不会自动形成层级。论文引入 pooling 或 wavelet，在若干层之间降低时间分辨率，使更高层以较少 positions 聚合更长时间范围。

\lecturefigure{slide-40-wavelet-motivation.jpg}{Wavelet motivation：在 intermediate embeddings 上显式构建多尺度层级。}{Stanford Online 官方视频 00:38:43--00:39:28。}

\begin{knowledgebox}{背景概念：Receptive field 与 Resolution}
Receptive field 指一个高层表示能够依赖的输入时间范围；resolution 指时间轴还保留多少位置。Pooling 扩大每个位置覆盖范围但丢失细粒度；wavelet 同时产生低频 approximation 与高频 detail，理论上能比纯平均保留更多变化，但如何组合仍决定信息是否真正被使用。
\end{knowledgebox}

\teachervoice{讲者把 wavelet 解释成“在一层里学习树状 hierarchy”。这不是说 wavelet 参数由模型学出；课堂使用的是固定 Haar-like operation，学习发生在 Transformer embeddings 与后续层。}

\subsection{Haar Wavelet：平均与差分}

对相邻 embeddings $h_{2i},h_{2i+1}\in\mathbb{R}^d$，最简单 Haar transform 可写成
\[
a_i=\frac{h_{2i}+h_{2i+1}}{\sqrt{2}},
\qquad
d_i=\frac{h_{2i}-h_{2i+1}}{\sqrt{2}}.
\]
$a_i$ 是低频 approximation，概括局部平均；$d_i$ 是 detail，保留相邻变化。若系统只保留平均，则等价于降采样 pooling；若保留或组合 detail，则能携带瞬态差异。长度从 $L$ 变为约 $L/2$ 后，后续 attention pair 数从 $L^2$ 降为约 $L^2/4$。

\lecturefigure{slide-41-wavelet-in-transformer.jpg}{在 Transformer layers 之间插入 Haar wavelet/average pathway。}{Stanford Online 官方视频 00:39:29--00:41:18。}

\begin{importantbox}{为什么“无新增参数”仍会改变模型}
Haar transform 没有 learned weights，但它改变 sequence length、频率响应、后续 attention graph 与 gradient path。参数量不变不等于计算、信息或 inductive bias 不变。
\end{importantbox}

\subsection{Architecture：Front End、Transformer 与 Temporal Reduction}

模型先把 waveform patches 输入 learned feed-forward front end，压成较低维 embeddings；随后通过多层 Transformer，并在中间使用 average pooling 或 wavelet pathway；最终聚合并输出多标签 scores。论文的 large model 使用 6 layers、64-dimensional embeddings、8 heads 和 feed-forward blocks，参数量约 2.3M，重点是与同表卷积 baselines 比较表示效率。

\lecturefigure{slide-42-audio-transformer-architecture.jpg}{Raw-waveform Audio Transformer：front-end patches、attention layers、temporal reduction 与 classifier。}{Stanford Online 官方视频 00:41:19--00:42:21。}

\lecturefigure{slide-43-learned-front-end.jpg}{Methodology：front end 学时频表示，Transformer 学 embedding，pooling/wavelet 扩展尺度。}{Stanford Online 官方视频 00:42:22--00:44:01。}

\begin{knowledgebox}{读图：不要把 Front End 当作“没有卷积就没有预处理”}
左侧仍把 waveform 组织为 patches，并用 2048-unit front end 与 64-dimensional projection 产生 tokens；Transformer 接收的是 learned embeddings。论文所谓无 convolution，指没有传统 convolutional backbone，不是输入 samples 不经过任何局部映射。
\end{knowledgebox}

\subsection{结果表：参数效率与指标范围}

Slide/论文表格显示：VGG-like 为 0.434 mAP，small Transformer 为 0.469，large 6-layer Transformer 为 0.525，加入 pooling 的同规模模型为 0.537。CRNN、ResNet-18 与 DenseNet-121 在该表中分别为 0.417、0.373、0.425。结果支持 raw-waveform Transformer 在该 FSD50K 设置中有竞争力，也提示 temporal reduction 有帮助。

\lecturefigure{slide-44-results-table.jpg}{Audio understanding 结果：卷积 baselines、Transformer 与 pooling/multi-scale variants。}{Stanford Online 官方视频 00:44:02--00:44:35；Audio Transformers v1 表 1。}

\begin{table}[H]
\centering
\small
\caption{FSD50K classroom-era comparison（mAP；参数量来自 source table）。}
\begin{tabular}{lrr}
\toprule
模型 & mAP & Parameters \\
\midrule
CRNN & 0.417 & 0.96M \\
VGG-like & 0.434 & 0.27M \\
ResNet-18 & 0.373 & 11.3M \\
DenseNet-121 & 0.425 & 12.5M \\
Small Transformer & 0.469 & 0.9M \\
Large 6-layer Transformer & 0.525 & 2.3M \\
Large 6-layer Transformer with pooling & 0.537 & 2.3M \\
\bottomrule
\end{tabular}
\end{table}

\begin{warningbox}{“Adieu Convolutions” 的合理结论}
表格说明在该数据、训练 recipe 与参数范围内，Transformer variants 达到更高 mAP；它没有控制所有现代 convolutional baselines，也没有比较 streaming latency、energy、robustness 或更大数据预训练。标题是研究方向宣言，不是普遍定理。
\end{warningbox}

\teachervoice{讲者对 pooling/wavelet 的提升感到意外，因为新增 learned parameters 很少。应把这理解为 inductive-bias evidence：改变时间层级本身可能比单纯加参数更重要。}

\subsection{本章小结}

Raw-waveform understanding 不等于把所有 samples 全局互联。Learned front end 先形成 tokens，pooling/wavelet 再降低时间分辨率。结果支持在 FSD50K 上使用 Transformer 与 multi-scale bias，但结论仍受数据、指标和 baseline 范围约束。

